% ============================================================================
%  Chapitre 3 — État de l'art
%  Cible : 12 à 15 pages. À rédiger J4-J5.
%
%  RÈGLE : chaque section se termine par un paragraphe « ce que l'on retient
%  pour ce travail ». Un état de l'art qui ne redescend pas sur le cas traité
%  est systématiquement sanctionné.
% ============================================================================
\chapter{État de l'art}
\label{chap:etat-art}

\section{Intégration et livraison continues}
\label{sec:art-cicd}

\subsection{Origines et définitions}
L'intégration continue est définie comme la pratique consistant à intégrer
fréquemment le travail de chaque développeur dans une ligne principale
partagée, chaque intégration étant validée par une construction et une
exécution automatisées des tests~\cite{fowler2006ci,duvall2007ci}. La livraison
continue étend ce principe en maintenant en permanence le logiciel dans un état
déployable~\cite{humble2010cd}.

\aecrire{Distinguer clairement intégration continue, livraison continue et
déploiement continu. Préciser laquelle de ces trois pratiques est visée dans ce
travail et pourquoi le déploiement continu n'est pas applicable à un produit
matériel certifié.}

\subsection{Bénéfices et coûts rapportés}
\aecrire{S'appuyer sur les études empiriques : Hilton et al. sur l'usage et le
coût de l'intégration continue~\cite{hilton2016usage}, Vasilescu et al. sur
les effets qualité et productivité~\cite{vasilescu2015quality}, Ståhl et Bosch
sur la variabilité des mises en œuvre industrielles~\cite{stahl2014modeling}.
Mentionner les coûts : maintenance de la chaîne, temps d'attente, tests
instables.}

\subsection{Revues systématiques}
La revue systématique de Shahin et al.~\cite{shahin2017continuous} recense les
approches, outils et difficultés associés à l'intégration, la livraison et le
déploiement continus, et identifie notamment la durée des constructions et la
complexité des environnements de test comme obstacles récurrents.

\subsection{Ce que l'on retient}
\aecrire{2 à 3 phrases reliant explicitement au cas traité.}

\section{Spécificités des systèmes embarqués}
\label{sec:art-embarque}

\aecrire{Cœur de l'originalité du sujet. Points à traiter : compilation
croisée, dépendance au matériel pour les tests, durée des constructions de
distributions complètes, cycles de certification, impossibilité de déployer en
continu. S'appuyer sur Mårtensson et al.~\cite{martensson2016embedded} et
Lwakatare et al.~\cite{lwakatare2019devops}.}

\subsection{Construction de distributions Linux embarquées}
\aecrire{Présenter Yocto/OpenEmbedded : notion de couche, de recette, de
machine, de distribution ; rôle du cache d'état partagé et du cache de
téléchargement dans la réduction du temps de construction. Présenter kas comme
outil de description déclarative de la configuration de construction. Citer la
documentation du projet Yocto~\cite{yoctoref}.}

\subsection{Ce que l'on retient}
\aecrire{Justifier la séparation en deux chaînes distinctes, applicative et
distribution, qui sera reprise au chapitre~4.}

\section{Reproductibilité de la construction}
\label{sec:art-reproductibilite}

La reproductibilité des constructions, c'est-à-dire la capacité à obtenir un
artefact binaire identique à partir d'une même révision du code source, est
présentée par Lamb et Zacchiroli comme un levier d'intégrité de la chaîne
d'approvisionnement logicielle~\cite{lamb2022reproducible}.

\aecrire{Articuler deux niveaux : (1) reproductibilité \emph{stricte} au sens
bit à bit, difficile à atteindre ; (2) reproductibilité \emph{de
l'environnement}, par conteneurisation et épinglage des versions, qui est
l'objectif réaliste retenu ici. Assumer ce choix explicitement, c'est un point
de discussion probable en soutenance.}

\section{Analyse statique de code}
\label{sec:art-sast}

\aecrire{Présenter le principe de l'analyse statique, puis les résultats
empiriques : Beller et al. sur l'état de l'analyse statique dans les projets
open source~\cite{beller2016analyzing}, Zampetti et al. sur l'intégration des
outils d'analyse statique dans les chaînes d'intégration
continue~\cite{zampetti2017static}. Traiter la question centrale des faux
positifs et de la stratégie de traitement de la dette initiale.}

\subsection{Politique de blocage}
\aecrire{Question de conception importante : faire échouer la construction
sur détection de défaut, ou seulement notifier ? Présenter les deux postures
et leurs conséquences sur l'adoption par les équipes. Ce point sera tranché
au chapitre~4 et discuté au chapitre~6.}

\section{Sécurité de la chaîne d'approvisionnement logicielle}
\label{sec:art-supply-chain}

\subsection{Menaces}
Ohm et al. proposent une taxonomie des attaques visant la chaîne
d'approvisionnement des logiciels libres~\cite{ohm2020backstabber}.

\subsection{Nomenclature logicielle}
La nomenclature logicielle (\gls{sbom}) recense les composants entrant dans la
composition d'un produit. Les éléments minimaux attendus ont été définis par la
NTIA~\cite{ntia2021sbom}, et le format SPDX a été normalisé sous la référence
ISO/IEC~5962~\cite{iso5962}.

\subsection{Analyse de composition logicielle}
\aecrire{Principe de l'analyse de composants binaires : identification des
composants embarqués dans un artefact compilé, corrélation avec les bases
\gls{cve}, scores \gls{cvss}. Expliquer pourquoi l'analyse binaire est plus
adaptée qu'une analyse des seuls fichiers de dépendances dans un contexte
Yocto.}

\section{Cadre normatif et réglementaire}
\label{sec:art-normes}

La norme IEC~62443-4-1~\cite{iec62443_4_1} définit les exigences de cycle de
développement sécurisé pour les produits d'automatisation industrielle. Le
règlement (UE) 2024/2847~\cite{cra2024} introduit des obligations de
cybersécurité pour les produits comportant des éléments numériques, dont la
tenue d'une nomenclature logicielle et le traitement des vulnérabilités
pendant la durée de support.

\aecrire{Établir explicitement la correspondance entre les pratiques mises en
œuvre au chapitre~5 et les exigences normatives. C'est un argument fort, et
c'est ce qui différencie ce mémoire d'un simple travail d'outillage.}

\aecrire{Compléter le tableau ci-dessous à partir du texte de la norme :
pratiques de gestion des défauts de sécurité, de vérification et de gestion des
composants tiers. Une ligne par exigence, avec renvoi vers la section du
chapitre~5 qui la couvre.}

\begin{table}[H]
\centering
\caption{Correspondance entre exigences normatives et pratiques mises en œuvre}
\label{tab:art-normes}
\begin{tabularx}{\textwidth}{@{}l X l@{}}
\toprule
\textbf{Exigence} & \textbf{Intitulé} & \textbf{Mise en œuvre} \\
\midrule
 & & \\
\bottomrule
\end{tabularx}
\end{table}

\section{Infrastructure as code et exécuteurs éphémères}
\label{sec:art-iac}

\aecrire{Présenter le principe de l'infrastructure décrite en code, à partir
de Morris~\cite{morris2020iac} et de la cartographie systématique de Rahman et
al.~\cite{rahman2019iac}. Traiter ensuite la question spécifique de la
sécurité des exécuteurs auto-hébergés : isolation entre exécutions,
non-persistance, gestion des secrets, portée des identités d'accès au
fournisseur cloud.}

\section{Mesure de la performance de livraison logicielle}
\label{sec:art-dora}

Les travaux du programme \gls{dora} identifient quatre indicateurs de
performance de livraison : la fréquence de déploiement, le délai de mise en
production, le taux d'échec des changements et le délai de
rétablissement~\cite{forsgren2018accelerate}.

\aecrire{Discuter honnêtement l'applicabilité de ces quatre indicateurs à un
produit embarqué qui n'est pas déployé en continu. Proposer les adaptations
retenues (par exemple : fréquence d'intégration au lieu de fréquence de
déploiement, délai de rétablissement d'une construction cassée). Cette
adaptation argumentée constitue en soi une contribution méthodologique.}

\section{Synthèse et positionnement}
\label{sec:art-synthese}

\aecrire{Une page. Reprendre les quatre sous-questions de recherche et
indiquer, pour chacune, ce que la littérature apporte et ce qui reste à
concevoir dans le cas industriel traité. Terminer par un tableau de synthèse
positionnant ce travail.}
